\clearpage
\begin{examplebox}
    \textbf{\textcolor{CExample}{例一：一次函数，直接代两端（最基础）}}
    \textbf{题目：}求 $y=2x+1$ 在 $[0,3]$ 上的值域。
    \begin{description}[style=nextline,leftmargin=2.8em,labelsep=0.5em,itemsep=0.35em]
        \item[\textbf{第一步：判断方向。}]
            $x$ 每增大 $1$，$y$ 就增大 $2$；函数一直递增，且图像是一条不断开的直线。
        \item[\textbf{第二步：算端点。}]
            $f(0)=1$，$f(3)=7$。
        \item[\textbf{第三步：写答案。}]
            \[\boxed{\text{值域}=[1,7]}.\]
    \end{description}
    \textbf{提速提醒：}这类题不需要列许多函数值，只需代入 $0$ 和 $3$。

    \medskip
    \textbf{\textcolor{CExample}{例二：递减函数，小的函数值反而在右边}}
    \textbf{题目：}求 $f(x)=5-2x$ 在 $[-1,2]$ 上的值域。
    \begin{description}[style=nextline,leftmargin=2.8em,labelsep=0.5em,itemsep=0.35em]
        \item[\textbf{第一步：判断方向。}]
            $x$ 增大时 $-2x$ 变小，所以函数一直递减且连续。
        \item[\textbf{第二步：算端点。}]
            $f(-1)=7$，$f(2)=1$。
        \item[\textbf{第三步：先小后大。}]
            \[\boxed{\text{值域}=[1,7]}.\]
    \end{description}
    \textbf{提速提醒：}写值域时不能写成 $[7,1]$。

    \medskip
    \begin{minipage}{\linewidth}
\textbf{\textcolor{CExample}{例三：二次函数，中途拐弯就不能只看两端}}
\begin{center}\begin{tikzpicture}[font=\small,>=stealth,line join=round]
\draw[gray!45,rounded corners=3pt] (0,0) rectangle (14.2,9.4);
\node at (7.1,8.8) {二次函数中途改变方向，先在顶点处分段};

\begin{scope}[shift={(2.2,2)},x=1.15cm,y=.46cm]
\draw[->] (-.3,0)--(4.8,0) node[right]{$x$};\draw[->] (0,-.2)--(0,12.7) node[above]{$y$};
\draw[orange!90!black,very thick,domain=0:1] plot(\x,{(\x-1)^2+2});\draw[blue,very thick,domain=1:4] plot(\x,{(\x-1)^2+2});
\foreach \x/\y in {0/3,1/2,4/11}{\fill (\x,\y) circle(2pt);}
\draw[dashed] (1,0)--(1,2);\draw[dashed] (4,0)--(4,11)--(0,11);\node[below] at (1,0){$1$};\node[below] at (4,0){$4$};\node[left] at (0,11){$11$};\node[below right] at (1,2){$(1,2)$};\node[above left] at (0,3){$3$};
\end{scope}
\node[align=left,anchor=west] at (8.3,5.1){$f(x)=(x-1)^2+2$\\[10pt] $[0,1]$ 递减：$[2,3]$\\[10pt] $[1,4]$ 递增：$[2,11]$\\[10pt] 合并值域 $[2,11]$};
\node[text width=13.2cm,align=center] at (7.1,.65) {若只看两个端点，会漏掉顶点附近低于 $3$ 的函数值。};
\end{tikzpicture}
\end{center}
\end{minipage}


    \textbf{题目：}求 $f(x)=x^2-2x+3$ 在 $[0,4]$ 上的值域。
    \begin{description}[style=nextline,leftmargin=2.8em,labelsep=0.5em,itemsep=0.4em]
        \item[\textbf{第一步：配方找最低点。}]
            \[
                f(x)=(x-1)^2+2.
            \]
            当 $x=1$ 时最低，为 $2$；这个点在 $[0,4]$ 里面。
        \item[\textbf{第二步：把区间拆成两段。}]
            $[0,1]$ 上递减：$f(0)=3$，$f(1)=2$，值域 $[2,3]$。
            $[1,4]$ 上递增：$f(1)=2$，$f(4)=11$，值域 $[2,11]$。
        \item[\textbf{第三步：把两段答案合起来。}]
            \[
                [2,3]\cup[2,11]=\boxed{[2,11]}.
            \]
    \end{description}
    \textbf{提速提醒：}如果只比较 $f(0)=3$ 和 $f(4)=11$，
    会把 $2$ 到 $3$ 之间的值漏掉。

    \medskip
    \textbf{\textcolor{CExample}{例四：绝对值函数，先拆再合并}}
    \textbf{题目：}求 $f(x)=|x-1|$ 在 $[-1,2]$ 上的值域。
    \begin{description}[style=nextline,leftmargin=2.8em,labelsep=0.5em,itemsep=0.4em]
        \item[\textbf{第一步：看 $x=1$ 两边。}]
            \[
                |x-1|=\begin{cases}
                1-x,&x\le1,\\
                x-1,&x\ge1.
                \end{cases}
            \]
        \item[\textbf{第二步：分别求值域。}]
            $[-1,1]$ 上由 $2$ 降到 $0$，值域 $[0,2]$；
            $[1,2]$ 上由 $0$ 升到 $1$，值域 $[0,1]$。
        \item[\textbf{第三步：合并。}]
            \[
                \boxed{[0,2]}.
            \]
    \end{description}
    \textbf{提速提醒：}绝对值函数最容易在绝对值里面等于零的位置改变方向。

    \medskip
    \textbf{\textcolor{CExample}{例五：带字母，判断上升还是下降再写值域}}
    \textbf{题目：}求 $f(x)=(m-2)x+1$ 在 $[1,3]$ 上的值域（$m$ 为实数）。
    \begin{description}[style=nextline,leftmargin=2.8em,labelsep=0.5em,itemsep=0.4em]
        \item[\textbf{第一步：端点值先算好。}]
            \[
                f(1)=m-1,\qquad f(3)=3m-5.
            \]
        \item[\textbf{第二步：比较两个端点。}]
            \[
                f(3)-f(1)=2(m-2).
            \]
            当 $m>2$ 时递增；$m<2$ 时递减；$m=2$ 时是常数函数。
        \item[\textbf{第三步：分三种情况写。}]
            \[
                \text{值域}=\begin{cases}
                [m-1,\ 3m-5],&m>2,\\
                [1,1],&m=2,\\
                [3m-5,\ m-1],&m<2.
                \end{cases}
            \]
    \end{description}
    \textbf{提速提醒：}一次函数题可以直接比较 $f(a)$ 与 $f(b)$ 的大小，
    不必重复计算许多函数值。
\end{examplebox}
